\begin{abstract}
Do language models have recurring ways of teaching that predict their behavior
on new problems? We investigate educational personality through observable
actions rather than self-reported human traits. An audit of a million-record
educational archive motivates a separate prospective study of five deployments,
four domains, controlled student cues, and repeated requests. We generate
5,320 tutor responses and lock behavior forecasts before confirmation
generation, using three model coders without new human annotation. On 32 new
source problems, model-default profiles reduce Brier loss beyond educational
context by 0.0913 for answer revelation, 0.1297 for reasoning elicitation, and
0.0198 for affect acknowledgement. Model-specific student-state patterns add
0.0097 beyond model-by-domain profiles for acknowledgement, with no resolved
gain for the other two actions. Explicit teaching instructions make demanded
actions converge, while permissible accompanying actions remain different.
Repetition, model-deletion, and neutral-wording analyses constrain stronger
claims of uniform or invariant traits. The findings support recurring,
prompt-dependent behavioral profiles of specified deployments, without inferring
inner psychology or learning gains. The pattern persists in content-screened
matched subsets, subject to limited agreement on natural errors.
\end{abstract}
